%
% 4-cas.tex
%
% (c) 2026 Prof Dr Andreas Müller
%
\section{Computeralgebra und die Renaissance der Algebra}
\kopfrechts{Computeralgebra und die Renaissance der Algebra}
Die auf präzisen Grundlagen aufbauende Analysis des 19.~und des
frühen 20.~Jahrhunderts hat nichts an den Rechenregeln geändert,
die für die Berechnung von Ableitungen und Integralen von Funktionen
verwendet werden.
\index{Computeralgebra}%
Dies ist Ausdruck des in
Abschnitt~\ref{buch:einleitung:fundament:subsection:funktionsbegriff}
angesprochenen veränderten Verständnisses einer Funktion als
Zuordnung statt als ``analytischer'' Ausdruck.

%
% Das Problem der Integration in geschlossener Form
%
\subsection{Das Problem der Integration in geschlossener Form}
Die Suche nach Antworten auf die Frage, welche Funktionen auch
eine ``analytische'' Darstellung haben, blieb auch im 19.~Jahrhundert
aktiv.
Dabei ist zunächst eine Definition des bisher immer in Anführungszeichen
gesetzten Begriffs des ``analytischen'' Ausdrucks nötig.
Es ist klar, dass zu diesen ``analytischen'' Funktionen mindestens
Polynome und rationale Funktionen einer Variablen gehören sollen.
Für diese Funktionen sind vollständige Methoden zur Berechnung von
Ableitungen und Stammfunktionen bekannt.
Bei letzteren treten aber unvermeidlich Logartihmusfunktionen auf.
\index{Logarithmusfunktion}%

Schon im Altertum wurden viele weitere Funktionen entdeckt, welche auch
für den Aufbau von ``analytischen'' Ausdrücken zugelassen werden sollten.
Dazu gehören zunächst die Wurzelfunktionen, auch Radikale genannt.
\index{Wurzelfunktion}%
\index{Radikal}%
Sie können als Lösungen von Polynomgleichungen definiert werden.
Solche Funktionen heissen auch algebraisch.
\index{algebraische Funktion}%
Ein bedeutende Entdeckung durch Abel und Galois im 19.~Jahrhundert
\index{Abel, Niels Henrik}%
\index{Galois, Evariste@Galois, Évariste}%
besagt, dass sich die Nullstellen von Polynomen mindestens fünften
Grades im allgemeinen nicht durch Radikale darstellen lassen.
Die Menge der algebraischen Funktionen ist also bereits grösser als
die Klasse der Wurzelausdrücke.
\index{Wurzelausdruck}%
Im Sinne einer einheitlichen algebraischen Theorie darf man erlauben,
alle algebraischen Funktionen zur Kathegorie der ``analytischen'' Ausdrücke
zu schlagen.

Gestattet man die Verwendung algebraischer Ausdrücke, kann man auch
das Auftreten komplexer Funktionen nicht mehr verhindern, denn
das Polynom $p(x) = x^2 + z^2$ definiert zum Beispiel über die Gleichung
$p(x)=0$ die algebraische Funktionen $\sqrt{-z^2}=\pm iz$.
Damit wird klar, dass die komplexe Analysis der richtige Ort ist, diese
Frage weiter zu untersuchen.
Aus diesem Grund werden komplexe Zahlen und komplexe Funktionen im
Kapitel~\ref{chapter:komplex} und ihre Ableitungen im
Kapitel~\ref{chapter:holomorph} studiert.
Es stellt sich heraus, dass Differenzierbarkeit eine sehr starke
Einschränkung ist.

Ebenfalls zu den ältesten Kandidatenfunktionen gehörten die
trigonometrischen Funktionen, die für die sphärische und
ebene Geometrie schon früh verwendet wurden.
\index{Geometrie, sphärische}%
\index{trigonometrische Funktion}%
Die eulersche Beobachtung $e^{ix}=\cos x + i\sin x$ bringt 
eine bedeutende Vereinfachung in diesen Zoo von neuen Funktionen.
Sie zeigt nämlich, dass sich alle diese Funktionen allein durch
die komplexe Exponentialfunktion $e^z$ beschreiben lassen.
\index{Exponentialfunktion}%
Auch die Umkehrfunktionen und die hyperbolischen Funktionen mit
\index{hyperbolische Funktionen}%
ihren Umkehrfunktionen lassen sich als Ausdrücke in $e^z$ und $\log z$
schreiben.

Die Exponential- und die Logarithmusfunktion sind jedoch nicht willkürlich.
Sie können beide durch ihr Verhalten unter Ableitung charakterisiert
werden, gilt doch
\begin{align}
\frac{d}{dx}e^{f(x)} &= e^{f(x)}\frac{df}{dx}(x)
&&\text{und}&
\frac{d}{dx}\log f(x) &= \frac{1}{f(x)}\frac{df}{dx}(x).
\label{buch:einleitung:cas:eqn:explog}
\end{align}
So wie algebraische Gleichungen Wurzeln von beliebigen rationalen
Funktionen zu definieren erlauben, könnnen die Eigenschaften
\eqref{buch:einleitung:cas:eqn:explog}
Exponential- und Logarithmusfunktionen beliebiger anderer Funktionen
$f(x)$ als Lösungen von Differentialgleichungen charakterisieren.

Schliesslich soll mit solchen Ausdrücken beliebig gerechnet werden
können.
Dies lässt sich algebraisch dadurch formulieren, dass die Menge der
zulässigen Funktionen einen Körper bilden soll.
\index{Körper}%
Dieser Körper enthält mindestens die rationalen Zahlen $\mathbb{Q}$,
\index{rationale Zahl}%
vielleicht noch weitere Konstanten wie $\pi$ oder $\gamma$ oder algebraische
Elemente wie $\!\sqrt{2}$ oder $\!\sqrt{\pi}$.
\index{pi@$\pi$}%
\index{gamma@$\gamma$}%
\index{2@$\sqrt{2}$}%
Er enthält aber auch die unabhängige Variable $x$, nach der abgeleitet
werden kann.
Ausserdem ist auf ihm die Differentiation als algebraische Operation
definiert, die den gewohnten Rechengesetzen der Ableitung genügt.

Zusammenfassend kann man also sagen, dass als Klasse von ``analytischen''
Ausdrücken, die die Theorie erfassen soll, alles verwendet werden kann,
was aus dem Grundkörper $\mathbb{Q}$ durch Körpererweiterung um
\index{Q@$\mathbb{Q}$}%
weitere Konstanten, algebraische Funktionen und Exponential- oder
Logarithmusfunktionen werden kann.
Solche Funktionen heissen \emph{elementare} Funktionen und werden
\index{elementare Funktion}%
im Abschnitt~\ref{buch:intalg:section:elementar} eingeführt.

Mit den elementaren Funktionen lässt sich jetzt das Integrationsproblem
exakt formulieren.

\begin{aufgabe}
\label{buch:einleitung:cas:aufgabe}
Zu einer elementaren Funktion $f(x)$ finde eine elementare Stammfunktion
$g(x)$ oder beweise, dass keine elementare Stammfunktion existiert.
\end{aufgabe}

Die Entscheidung, dass eine Funktion $f(x)$ keine elementare
Stammfunktion hat, bedeutet auch, dass es sinnvoll ist, das unbestimmte
Integral als ``neue'' Funktion zu definieren.
So wurde die Stammfunktion von $e^{-x^2}$, welche nicht elementar ist,
als die Fehlerfunktion definiert.
Der Beweis, dass $e^{-x^2}$ nicht elementar ist, wird in
Abschnitt~\ref{buch:intalg:section:risch} mit Hilfe des Risch-Algorithmus
gegeben.
Auch die elliptischen Integrale, die in der Arbeit von Baris Catan 
in Kapitel~\ref{chapter:elliptisch} eingeführt werden, sind nicht
elementare Stammfunktionen elementarer Funktionen.

Ein Funktionenkörper mit einer Derivation kennt das Konzept der
Zusammensetzung von Funktionen nicht und kann daher auch nicht
sinnvoll eine Kettenregel formulieren.
Die Definitionen \eqref{buch:einleitung:cas:eqn:explog} für
Exponential- und Logarithmusfunktionen führen die Kettenregel
durch die Hintertür wieder ein.
Etwas Ähnliches passiert auch für die Wurzelfunktion, die durch
$w(x)^2-f(x)=0$ definiert ist.
Die Lösung dieser Gleichung ist $w(x)=\!\sqrt{f(x)}$.
Für die Ableitung von $w(x)$ wird die Kettenregel nicht benötigt, 
da man sie durch Ableiten der definierenden Gleichung finden kann:
\[
0
=
\frac{d}{dx}(w(x)^2-f(x))
=
2w(x)\frac{dw}{dx}(x) - \frac{df}{dx}(x)
\qquad\Rightarrow\qquad
\frac{dw}{dx}(x) = \frac{1}{2w(x)}\frac{df}{dx}(x).
\]
Dies ist die Kettenregel für die Quadratwurzelfunktion.
Für einen Funktionenkörper aus elementaren Funktionen sind damit
alle nötigen Versionen der Kettenregel automatisch erfüllt, die
Kettenregel wie auch das Konzept der Zusammensetzung von Funktionen
sind daher für die Lösung der Aufgabe~\ref{buch:einleitung:cas:aufgabe}
entbehrlich.

%
% Das Liouville-Prinzip
%
\subsection{Das Liouville-Prinzip}
Eine erste systematische Untersuchung des Problems der Integration
in geschlossener Form stammt von Joseph Liouville.
Einige seiner Resultate waren aber schon früher bekannt.
Liouville fand eine vollständige Charakterisierung der elementaren
\index{Liouville, Joseph}%
Funktionen, die auch eine elementare Stammfunktion haben.
Hat man bereits einen Körper von Funktionen konstruiert, in dem sich
der Integrand befindet, dann haben genau diejenigen Funktionen eine
elementare Stammfunktion, die die Form
\begin{align*}
f(x)
&=
v'_0(x)
+
\sum_{i=1}^n c_i \frac{v_i'(x)}{v_i(x)}
\intertext{haben, wobei $c_i$ Konstanten sind und $v_0(x)$ und $v_i(x)$
Funktionen im konstruierten Körper sind.
Die Stammfunktion hat dann die Form}
\int f(x)\,dx
&=
v_0(x)
+
\sum_{i=1}^c c_i\log v_i(x)
\end{align*}
haben.
\index{Liouville-Prinzip}%
Wenn ein Integrand eine elementare Stammfunktion hat, dann muss
Sie eine Linearkombination von den bereits konstruierten Funktionen
sowie möglicherweise von neuen Logarithmusfunktionen sein.

%
% Computeralgebra und der Algorithmus von Risch
%
\subsection{Computeralgebra und der Algorithmus von Risch}
Mit der Entwicklung ausreichend leistungsfähiger Computer in der zweiten
Hälfte des 20.~Jahrhunderts eröffnete sich die Möglichkeit, eine
Stammfunktion in geschlossener Form maschinell zu finden.
Das Liouville-Prinzip macht zwar eine Aussage über die Existenz einer
Stammfunktion in geschlossener Form, um diese zu berechnen, ist jedoch
mehr nötig.
Es muss ein Algorithmus formuliert werden, der zu einer elementaren
Funktion entweder die elementare Stammfunktion findet, falls sie
existiert, oder anderen Fall einen Beweis für ihre Nichtexistenz
konstruiert.
Ein solcher Algorithmus ist heute unter dem Namen Risch-Algorithmus
bekannt.
Er wird in Abschnitt~\ref{buch:intalg:section:risch} dargestellt.

Der Risch-Algorithmus hat nichts mehr mit der Analysis zu tun, die
Cauchy oder Weierstrass entwickelt haben.
\index{Risch-Algorithmus}%
Er beschreibt eine rein algebraische Vorgehensweise, um die Stammfunktion
zu finden.
Die bei numerischen Rechnung mit dem Computer allgegenwärtigen
Gleitkommazahlen kommen nur vor, wenn der Benutzer so wagemutig ist,
\index{Gleitkommazahl}%
solche in der Problemstellung explizit vorzugeben.
$\varepsilon$ und $\delta$ sind in dieser Welt abwesend und die
\index{epsilon@$\varepsilon$}%
\index{delta@$\delta$}%
die gefundenen Resultate sind nicht nur Annäherungen,
sondern exakt.

%
% Erweiterte Methoden
%
\subsection{Erweiterte Methoden}
Der Risch-Algorithmus in seiner Grundform ist limitiert auf elementare
Funktionen, also auf Funktionen, die sich mit Exponentialfunktionen,
Logarithmen und Wurzelfunktionen aus rationalen Funktionen konstruieren
lassen.

Der Risch-Algorithmus arbeitet rekursiv.
Er zerlegt ein Integral in einzelne Summanden, die einfacher auszuwerten
sind und die auf die einzelnen Terme in der vom Liouville-Prinzip vorgegebenen
Linearkombination führen.
Ein solcher rekursiver Aufruf des Risch-Algorithmus kann zu einem
Integral führen, welches nicht als elementare Funktion geschrieben
werden kann.
Der Algorithmus muss sich damit aber nicht begnügen.
Viele berühmte Funktionen haben keine elementare Stammfunktion, zum
Beispiel
\begin{equation}
\frac{\sin x}{x},
\quad
\frac{1}{\log x},
\quad
e^{-x^2/2}
\quad\text{oder}
\quad
\frac{1}{\sqrt{(1-x^2)(1-k^2x^2)}}.
\label{buch:einleitung:cas:eqn:integranden}
\end{equation}
Es wurden daher neue spezielle Funktionen als Stammfunktionen dieser
Integranden definiert.
Der \emph{Integralsinus} ist die Stammfunktion
\index{Integralsinus}%
\begin{align*}
\operatorname{Si}(x)
&=
\int_0^x \frac{\sin t}{t}\,dt,
\intertext{der \emph{Integrallogarithmus} ist
\index{Integrallogarithmus}}
\operatorname{li}(x)
&=
\int_0^x \frac{dt}{\log t},
\intertext{die \emph{Fehlerfunktion} ist
\index{Fehlerfunktion}}
\operatorname{erf}(x)
&=
\frac{2}{\!\sqrt{\pi}} 
\int_0^x e^{-t^2}\,dt.
\intertext{Der letzte Integrand in \eqref{buch:einleitung:cas:eqn:integranden}
führt auf das unvollständige elliptische Integral}
F_J(x;k)
&=
\int_0^x \frac{dt}{\!\sqrt{(1-t^2)(1-k^2t^2)}}
\end{align*}
\index{elliptisches Integral}%
\index{Integral, elliptisch}%
erster Art in Jacobi-Form.
Moderne Computeralgebrasysteme verstehen diese speziellen Funktionen
und sind damit in der Lage, auch solche Integrale auszuwerten
und gegebenenfalls in Stammfunktionen auch mit solchen wohlbekannten
Funktionen darzustellen.

Eine weitere, in diesem Buch nicht weiter untersuchte Möglichkeit ist,
Funktionen als hypergeometrische Funktionen zu schreiben.
\index{hypergeometrische Funktion}%
\index{Funktion!hypergeometrisch}%
Diese sind durch eine Potenzreihe
\[
f(z)
=
\sum_{k=0}^\infty \frac{a_k}{k!}z^k
\]
definiert, deren aufeinanderfolgende Koeffizienten $a_k$ und $a_{k+1}$
Quotienten haben, die eine rationale Funktion von $k$ sind.
Für solche Funktionen ist ebenfalls eine reichhaltige Theorie 
verfügbar, die zum Beispiel in \cite{buch:seminarspeziellefunktionen}
entwickelt wird.
Die Ableitung oder Stammfunktion einer solchen hypergeometrischen Funktion
ist wiederum eine hypergeometrische Funktion.
Wenn es gelingt, einen Integranden als eine hypergeometrische Funktion
zu identifizieren, dann kann auf diesem Weg auch ein Ausdruck der
Stammfunktion als hypergeometrische Funktion gefunden werden.

Die hypergeometrischen Funktionen sind zwar nicht so gut bekannt wie
die elementaren Funktionen, aber viele von Ihnen werden in
Programmbibliotheken für das wissenschaftliche Rechnen in hoher Qualität
implementiert.
Die Darstellung der Stammfunktion als hypergeometrische Funktion
ermöglicht daher, für die numerische Berechnung numerische
Integrationsalgorithmen vermeiden, die fast immer von minderer
Genauigkeit wären.

%
% Algebra und Analysis
%
\subsection{Algebra und Analysis}
Die historische Entwicklung zeigt also, wie sich aus dem erst vor
allem für die symbolische Rechnung entwickelten Kalkül durch
Präzisierung der Grundlagen zunächst die moderne Analsysis mit
ihrer $\varepsilon$-$\delta$-Technik entwickelt hat.
Die algebraischen Aspekte der Operationen des Differenzierens
und Integrierens rückten dabei immer mehr in den Hintergrund und
wurden durch numerische Methoden ersetzt, die die Errungenschaften
der modernen rechnergestützten Wissenschaften ermöglich haben.

In der komplexen Analysis des 19.~Jahrhunderts, wie sie von Cauchy
\index{Analysis!komplexe}%
\index{komplexe Analysis}%
und anderen entwickelt wurde, war aber auch algebraische Struktur
angelegt, mit der Liouville seine Ideen zur Integrierbarkeit
elementarer Funktionen entwickeln konnte.
Lange Zeit trat diese in den Hintergrund, weil sich 
Wetterprognosen, Strömungssimulationen oder schwarze Löcher nicht
an die Einschränkungen halten, nur elementare Funktionen zu verwenden.

Moderne Computer haben der Algebra wieder den Platz eingeräumt, den
sie verdient.
Heute kann jeder Computer-Hobbyist mit einem Raspberry Pi für wenig
\index{Raspberry Pi}%
Geld einen Computer erstehen, auf dem ihm in Form der von Wolfram
Research für diese Platform gratis zur Verfügung gestellten Software
\emph{Mathematica} eine vollständige Implementation
\index{Mathematica}%
des Risch-Algorithmus zur Verfügung steht und ihm die symbolische Berechnung
der kompliziertesten Integrale ermöglicht, sofern eine solche mit den
bekannten Algorithmen möglich ist.

In diesem Buch wird daher zunächst die komplexe Analysis mit besonderem
Fokus auf die algebraischen Eigenschaften der Ableitung entwickelt.
In Kapitel~\ref{chapter:ratint} wird dann das Problem der Integration
rationaler Funktionen als algebraisches Problem adressiert und damit
die Grundlage für das Verständnis des Problems der Integration in
geschlossener Form gelegt.
Kapitel~\ref{chapter:intalg} zeigt dann, wie die algebraische
Struktur mit dem Risch-Algorithmus einen praktisch umsetzbaren
Integrationsalgorithmus liefern kann.



